% =====================================================================
%  Programación y automatización en LaTeX — suite de documentos LaTeX
%  Nivel 4 (Avanzado) · clase article.
%  Compilar: latexmk -pdf programacion.tex
% =====================================================================
\documentclass[11pt,a4paper]{article}

% --- Base estándar ---
\usepackage[utf8]{inputenc}
\usepackage[T1]{fontenc}
\usepackage{lmodern}
\usepackage{microtype}
\usepackage[spanish,es-noshorthands]{babel}
\usepackage{csquotes}
\usepackage{amsmath}\usepackage{amssymb}
\usepackage{booktabs}
\usepackage[margin=2.4cm]{geometry}
\usepackage{enumitem}
\usepackage{xcolor}
\usepackage{listings}
% --- Programación ---
\usepackage{ifthen}     % \ifthenelse
\usepackage{etoolbox}   % toggles, \ifdef...
\usepackage{tikz}       % trae \foreach (pgffor)
\usepackage{hyperref}
\hypersetup{unicode, pdftitle={Programación y automatización en LaTeX},
            pdfauthor={Rodrigo Martín Vidal Galán},
            colorlinks=true, linkcolor=blue!50!black, urlcolor=blue!50!black}
\usepackage[spanish,capitalise,nameinlink]{cleveref}

% --- Estilo de código ---
\definecolor{codeback}{RGB}{245,247,250}
\definecolor{codekw}{RGB}{20,60,120}
\definecolor{codecom}{RGB}{110,120,135}
\lstset{
  basicstyle=\ttfamily\footnotesize,
  keywordstyle=\color{codekw}\bfseries,
  commentstyle=\color{codecom}\itshape,
  backgroundcolor=\color{codeback},
  frame=single, rulecolor=\color{codekw!30},
  breaklines=true, showstringspaces=false, columns=fullflexible,
  language=[LaTeX]TeX, xleftmargin=8pt, framesep=4pt, extendedchars=true,
  literate={á}{{\'a}}1 {é}{{\'e}}1 {í}{{\'i}}1 {ó}{{\'o}}1 {ú}{{\'u}}1
           {ñ}{{\~n}}1 {¿}{{?`}}1 {¡}{{!`}}1,
}
\newcommand{\C}[1]{\texttt{\textbackslash #1}}
\newcommand{\pkg}[1]{\textsf{#1}}
\newcommand{\produce}{\par\smallskip\noindent\textit{produce:}\quad}

% --- Macros de demostración (se usan en el cuerpo) ---
\newcommand{\saludo}{¡Hola desde una macro!}
\newcommand{\cuadrado}[1]{#1^{2}}
\newcommand{\vect}[1][v]{\mathbf{#1}}      % argumento OPCIONAL (default: v)
\newcounter{ejemplo}                        % contador propio

\title{\textbf{Programación y automatización en \LaTeX}\\[2pt]
       \large Macros, lógica, bucles y generación de contenido}
\author{Rodrigo Martín Vidal Galán}
\date{Junio de 2026}

\begin{document}
\maketitle

\noindent \LaTeX{} no es solo un procesador de texto: es un \textbf{lenguaje de
programación} (de hecho, Turing-completo). Eso permite definir tus propios
comandos, automatizar lo repetitivo, tomar decisiones lógicas, repetir con
bucles e incluso \textbf{inyectar} resultados de otros lenguajes. Acá va lo
esencial, de lo cotidiano (macros) a lo avanzado (\pkg{expl3}, Lua, generación).

\newpage
\tableofcontents
\newpage

% =====================================================================
\section{Macros: tus propios comandos}
% =====================================================================
La herramienta más usada. \C{newcommand} define un comando nuevo; si ya existe,
\LaTeX{} avisa (a diferencia de \C{def}, que pisa en silencio):
\begin{lstlisting}
\newcommand{\saludo}{Hola desde una macro!}     % sin argumentos
\newcommand{\cuadrado}[1]{#1^{2}}               % 1 argumento: #1
\newcommand{\vect}[1][v]{\mathbf{#1}}           % argumento OPCIONAL (default v)
\end{lstlisting}
\produce \saludo{} \quad $\cuadrado{x}$, \quad $\cuadrado{(a+b)}$, \quad
$\vect$ y $\vect[w]$.
\medskip

\begin{itemize}[noitemsep,leftmargin=*]
  \item \texttt{[n]} declara \textbf{n} argumentos, accesibles como
        \texttt{\#1\dots\#9}.
  \item \texttt{[n][def]} hace que el \textbf{primer} argumento sea opcional, con
        valor por defecto \texttt{def}.
  \item \C{renewcommand} \textbf{redefine} uno existente; \C{providecommand} lo
        define solo si no existe.
\end{itemize}
\textbf{Por qué importa:} una macro centraliza una notación. Si definís
\C{newcommand\{\textbackslash R\}\{\textbackslash mathbb\{R\}\}}, escribís
\C{R} en todos lados y, si algún día querés otro estilo, lo cambiás \emph{en un
solo lugar} (principio DRY).

\paragraph{Entornos propios.} \C{newenvironment} hace lo mismo para pares
\texttt{begin/end}:
\begin{lstlisting}
\newenvironment{remarca}{\par\noindent\textbf{Nota:~}\itshape}{\par}
\end{lstlisting}

% =====================================================================
\section{Contadores y longitudes}
% =====================================================================
\LaTeX{} cuenta con \textbf{contadores} (enteros) y mide con \textbf{longitudes}.
Podés crear y manipular los tuyos:
\begin{lstlisting}
\newcounter{ejemplo}
\stepcounter{ejemplo}            % +1
\setcounter{ejemplo}{5}          % asignar
\arabic{ejemplo} / \roman{...}   % mostrarlo (1.. / i..)
\the\value{ejemplo}              % su valor

\newlength{\misangria}\setlength{\misangria}{1.5em}
\addtolength{\misangria}{2pt}
\end{lstlisting}
\stepcounter{ejemplo}\stepcounter{ejemplo}\stepcounter{ejemplo}%
\produce el contador vale \arabic{ejemplo} (en romano: \roman{ejemplo}).
\medskip

Los contadores son la base de la numeración automática (secciones, ecuaciones,
teoremas): todos son contadores que \LaTeX{} incrementa solo.

% =====================================================================
\section{Lógica: condicionales}
% =====================================================================
Para que el documento \textbf{decida}, hay varias vías. La más legible es
\C{ifthenelse} (paquete \pkg{ifthen}):
\begin{lstlisting}
\ifthenelse{\value{ejemplo} > 3}{es grande}{es chico}
\ifthenelse{\equal{\modo}{borrador}}{...}{...}
\end{lstlisting}
\produce el contador \arabic{ejemplo}
\ifthenelse{\value{ejemplo} > 3}{ \textbf{es grande}}{ es chico}.
\medskip

Con \pkg{etoolbox} se trabaja con \textbf{interruptores} (\emph{toggles}), muy
cómodos para activar/desactivar partes:
\begin{lstlisting}
\newtoggle{soluciones}
\toggletrue{soluciones}     % o \togglefalse
\iftoggle{soluciones}{\textbf{Solucion:} ...}{}
\end{lstlisting}
También existen los condicionales primitivos \C{ifnum}, \C{ifdim}, y
\C{ifdefined} para comprobar si un comando existe.

% =====================================================================
\section{Bucles}
% =====================================================================
Para repetir, lo más práctico es \C{foreach} (de \pkg{pgffor}, que trae
\pkg{tikz}): itera sobre una lista o un rango:
\begin{lstlisting}
\foreach \n in {1,2,3,4,5}{\n\ }          % lista explicita
\foreach \n in {1,...,10}{$\cuadrado{\n}=\the\numexpr\n*\n\ $ }  % rango
\end{lstlisting}
\produce \foreach \n in {1,...,6}{$\n^2=\the\numexpr\n*\n\relax$\quad}
\medskip

\C{foreach} es ideal para generar tablas, grillas o figuras repetitivas sin
copiar y pegar. También existen \C{loop\dots\C{repeat}} (TeX) y los bucles de
\pkg{expl3}.

\paragraph{Ejemplo: una fila de puntos con \texttt{\textbackslash foreach}.}
\begin{center}
\begin{tikzpicture}
  \foreach \x in {0,1,...,8}{ \fill[blue!60] (\x*0.6,0) circle (2.5pt); }
\end{tikzpicture}
\end{center}

% =====================================================================
\section{\texttt{expl3}: la capa de programación moderna (LaTeX3)}
% =====================================================================
La programación \enquote{clásica} de TeX es poderosa pero críptica. \pkg{expl3}
(ya parte del \LaTeX{} actual) ofrece una sintaxis \textbf{ordenada}, con
funciones y tipos de datos (listas, secuencias, \emph{property lists}), pensada
para escribir paquetes:
\begin{lstlisting}
\ExplSyntaxOn
\cs_new:Npn \mi_doble:n #1 { \int_eval:n { 2 * #1 } }
\ExplSyntaxOff
\end{lstlisting}
Convenciones: los nombres llevan \texttt{\_} y un sufijo de tipo (\texttt{:n},
\texttt{:N}\dots), y dentro de \texttt{\textbackslash ExplSyntaxOn} el espacio
\textbf{no cuenta} (se usa \texttt{\textasciitilde} para espacio explícito), lo
que hace el código legible. Es el estándar para desarrollo serio de paquetes; para
el día a día alcanza con \C{newcommand}.

% =====================================================================
\section{Inyectar otros lenguajes}
% =====================================================================
\paragraph{Lua (con LuaLaTeX).} El motor \texttt{lualatex} incluye el lenguaje
\textbf{Lua}; podés ejecutar código durante la compilación:
\begin{lstlisting}
\directlua{ tex.print(2^10) }        % imprime 1024
% o el entorno del paquete luacode para bloques mas largos
\end{lstlisting}

\paragraph{Python (con \pkg{pythontex}).} Permite \textbf{correr Python} y volcar
el resultado en el PDF (requiere \texttt{-shell-escape} y un paso extra
\texttt{pythontex}):
\begin{lstlisting}
\begin{pycode}
import numpy as np
x = np.linspace(0, 1, 5)
print(x.mean())
\end{pycode}
\py{2 + 2}        % inserta 4 en el texto
\end{lstlisting}
Es ideal para que cálculos, tablas o gráficos provengan de datos reales sin
copiar números a mano (\enquote{documento reproducible}).

\paragraph{Generar y \texttt{\textbackslash input}.} Un patrón general: un script
externo (Python, etc.) genera un \texttt{.tex} (una tabla, por ejemplo) y el
documento lo incorpora con \C{input\{tabla\_generada.tex\}}. Así separás el
cálculo de la presentación.

% =====================================================================
\section{Compilación condicional}
% =====================================================================
Producir \textbf{varias versiones} del mismo fuente (con/sin soluciones, borrador
vs.\ final) sin duplicar el archivo:
\begin{lstlisting}
% con un toggle de etoolbox (ver seccion de logica)
\newtoggle{soluciones}\toggletrue{soluciones}
... \iftoggle{soluciones}{\textbf{Sol:} 42}{} ...

% o con el paquete comment para excluir bloques enteros:
\usepackage{comment}
\excludecomment{soluciones}   % todo lo de ese entorno se omite
\begin{soluciones} ... \end{soluciones}
\end{lstlisting}
Para borrador vs.\ final también sirven la opción \texttt{draft} de la clase y
\C{includeonly} (ver \emph{Documentos largos}).

% =====================================================================
\section{Algoritmos y pseudocódigo}
% =====================================================================
Para escribir \textbf{pseudocódigo} con formato (palabras clave, indentación,
numeración de líneas) hay paquetes dedicados —\pkg{algorithm2e} o
\pkg{algpseudocodex}—, en vez de simular a mano:
\begin{lstlisting}
\begin{algorithm}
\caption{Euclides: mcd(a,b)}
\While{$b \neq 0$}{
  $t \gets b$\;
  $b \gets a \bmod b$\;
  $a \gets t$\;
}
\Return $a$\;
\end{algorithm}
\end{lstlisting}
(Para \textbf{mostrar} código fuente con resaltado —no pseudocódigo— se usan
\pkg{listings} o \pkg{minted}; ver \emph{Estándares}.)

% =====================================================================
\section{Buenas prácticas}
% =====================================================================
\begin{itemize}[noitemsep,leftmargin=*]
  \item \C{newcommand} (no \C{def}) para tus macros: avisa si pisás algo.
  \item Macros \textbf{semánticas} y centralizadas (DRY); preámbulo reutilizable.
  \item Lógica legible con \pkg{ifthen}/\pkg{etoolbox}; \C{foreach} para repetir.
  \item \pkg{expl3} solo si desarrollás paquetes; para escribir, alcanza lo básico.
  \item Datos reales $\to$ \pkg{pythontex}/Lua o generá \texttt{.tex} y
        \C{input} (reproducibilidad).
\end{itemize}

% =====================================================================
\section*{Ver también}
% =====================================================================
\addcontentsline{toc}{section}{Ver también}
\begin{itemize}[noitemsep,leftmargin=*]
  \item \emph{Estándares} — \C{newcommand} vs.\ \C{def}, mostrar código.
  \item \emph{Documentos largos} — \C{includeonly}, opción \texttt{draft}.
  \item \emph{Flujo e instalación} — \texttt{-shell-escape} para pythontex/minted.
  \item Manuales: \texttt{texdoc usrguide} (LaTeX), \texttt{texdoc expl3},
        \texttt{texdoc pythontex}, \texttt{texdoc etoolbox}.
\end{itemize}

\end{document}
